\section{A current obtained by separation}\label{sec:currents}

We first construct the current that would obstruct compatibility:
a positive current and an anti-invariant correction whose sum
annihilates every closed two-form.

Give \(X\) its complex orientation. Fix a smooth \(J\)-Hermitian
metric \(g\), and write
\[
F(u,v)=g(Ju,v),\qquad
(R\alpha)(u,v)=\tfrac12\bigl(\alpha(u,v)-\alpha(Ju,Jv)\bigr).
\]
Thus \(R\) is the orthogonal projection onto the anti-invariant
two-forms, and \(1-R\) projects onto the invariant two-forms.
All forms and currents are real.

We identify bivectors with two-forms using \(g\). Let
\(\mathscr C\to X\) be the bundle of unit complex line bivectors
\[
\mathscr C_x=\{v\wedge Jv:\ |v|_g=1\}.
\]
For \(\ell\in\mathscr C_x\), let \(L_\ell\subset T_xX\) be its
oriented plane. The following pointwise identities will be used:
\begin{equation}\label{eq:algebra}
\im R\subset\Lambda_g^+,\qquad
*F=F,\qquad |F|^2=2,\qquad
|\ell|=1,\qquad \ell^+=F/2,\qquad F(\ell)=1.
\end{equation}
They follow by taking a unitary frame
\((v,Jv,w,Jw)\). In that frame \(F\) corresponds to
\(v\wedge Jv+w\wedge Jw\), and the anti-invariant subspace is spanned
by
\[
v\wedge w-Jv\wedge Jw,\qquad
Jv\wedge w+v\wedge Jw.
\]

A two-current acts on smooth two-forms. It is \emph{closed} if it
annihilates exact two-forms. Under the \(L^2\) identification with
distributional two-forms, this means \(d^*T=0\), or equivalently
\(d(*T)=0\). Multiplication and orthogonal projections on
distributional forms act on currents by duality.

\begin{lemma}[A closed lift]\label{lem:lift}
There is a real order-zero pseudodifferential operator
\[
B:C^\infty(\im R)\longrightarrow\Omega^2(X)
\]
such that \(dB=0\) and \(RB=\Id\).
\end{lemma}

\begin{proof}
Consider \(\cL=Rdd^*\) on sections of \(\im R\). If \(\xi\ne0\),
the projection of
\(\xi\wedge\iota_{\xi^\sharp}\) from self-dual two-forms back to
self-dual two-forms is \(|\xi|^2/2\) times the identity.
By \eqref{eq:algebra}, the principal symbol of \(\cL\) on
\(\im R\) is therefore the same scalar. In particular \(\cL\)
is elliptic. For anti-invariant smooth forms \(a,b\),
\[
\ip{\cL a}{b}_{L^2}=\ip{d^*a}{d^*b}_{L^2}.
\]
It is nonnegative and self-adjoint. Every element of its kernel
is smooth and coclosed; being self-dual, it is also closed.

Let \(H\) be the orthogonal projection onto this kernel and let
\(G\) be the inverse of \(\cL\) on its orthogonal complement,
extended by zero on the kernel. The elliptic generalized-inverse
theorem gives \(G\) of order \(-2\) and \(H\) a smoothing
finite-rank operator; see
\cite[Proposition~6.4 and Theorem~6.3]{MelroseIML}.
Set
\[
B=dd^*G+H.
\]
Then \(dB=0\), while \(RB=\cL G+H=\Id\).
\end{proof}

The projected elliptic operator used here also appears in Lejmi's
proof of local compatibility in dimension four
\cite[proof of Theorem~1]{Lejmi2006}.
Dr\u{a}ghici and Zhang's exact-form result gives the corresponding
global linear correction phenomenon
\cite[Proposition~3.1]{DraghiciZhang2012}.
The present formula specifies the closed lift and its mapping properties;
positivity remains a separate problem.

\begin{proposition}\label{prop:separation}
If \(J\) admits no compatible symplectic form, there are currents
\(P,Q,T\) and a finite positive measure \(\lambda\) on
\(\mathscr C\) such that
\begin{align}
P(\alpha)&=\int_{\mathscr C}\alpha_x(\ell)\dd\lambda(x,\ell),
&\lambda(\mathscr C)&=1,\label{eq:positive-representation}\\
T&=P+Q,
&RP&=0,\quad RQ=Q,\quad *Q=Q,\label{eq:decomposition}\\
T(\alpha)&=0&&\text{for every smooth closed two-form }\alpha.
\label{eq:annihilation}
\end{align}
The current \(P\) annihilates every closed invariant two-form.
The projection \(\mu\) of \(\lambda\) to \(X\) is the
\emph{trace measure} of \(P\): for every smooth function \(b\),
\(P(bF)=\int_Xb\dd\mu\). Each of \(P,Q,T\) belongs to \(H^{-3}\).
\end{proposition}

\begin{proof}
In the Fr\'echet space of smooth invariant two-forms, the cone of
positive definite forms is open and convex. It is disjoint from
the linear subspace of closed invariant forms by hypothesis.
Hahn--Banach separation gives a nonzero continuous functional
nonnegative on the cone and zero on that subspace. Extend it to
all two-forms by composing with \(1-R\), and call it \(P\).
This is the positive-current separation framework of
Sullivan and Harvey--Lawson \cite{Sullivan1976,HarveyLawson1983};
for its almost-complex formulation, see
\cite[Section~3]{LiZhang2009}.

Positivity extends to semidefinite invariant forms by adding
\(\eps F\) and taking \(\eps\downarrow0\).
Since \(C\|\alpha\|_\infty F\pm(1-R)\alpha\) are semidefinite
for a fixed norm-comparison constant,
\[
|P(\alpha)|\le C\|\alpha\|_\infty P(F).
\]
Thus \(P\) has order zero, and \(P(F)>0\) because \(P\ne0\).
Scale \(P\) so that \(P(F)=1\).

For completeness, the measure representation can be obtained
in local unitary frames. Invariant two-forms identify with
Hermitian symmetric forms by
\(\alpha\mapsto\alpha(\,\cdot,J\,\cdot)\).
The positive functional has a positive Hermitian matrix of
coefficient measures. Its trace is a positive measure \(\mu\);
the coefficient measures are absolutely continuous with
respect to \(\mu\). Their Radon--Nikodym density is a positive
semidefinite Hermitian matrix of trace one, or equivalently
a convex combination of bivectors in \(\mathscr C_x\).
Measurable spectral decomposition in measurable unitary frames
gives \(\lambda\) in \eqref{eq:positive-representation}.
The identity \(F(\ell)=1\) shows that its projection is \(\mu\)
and that its total mass is one.

Define
\[
T(\alpha)=P(\alpha-BR\alpha),\qquad
Q(\alpha)=-P(BR\alpha).
\]
For closed \(\alpha\), the form \(\alpha-BR\alpha\) is closed
and invariant, proving \eqref{eq:annihilation}. Moreover \(RP=0\)
and \(RQ=Q\), so \(*Q=Q\) by \eqref{eq:algebra}.
In particular \(T\) is a closed current.

A finite measure on a compact four-manifold defines an element
of \(H^{-3}\), by the Sobolev embedding \(H^3\subset C^0\).
The order-zero operator \(RB^*\) is bounded on this Sobolev
space by the pseudodifferential Sobolev mapping theorem
\cite[Proposition~0.5.E]{Taylor1991}, so the same conclusion
holds for \(Q=-RB^*P\) and \(T\).
\end{proof}

Hereafter \(C\), with or without a subscript, denotes a finite
constant independent of the small scale parameters. It may
change from one occurrence to the next. Constants are uniform
over centers in \(X\).
